\documentclass[11pt]{article}
\usepackage[a4paper,margin=25mm]{geometry}
\usepackage[T1]{fontenc}
\usepackage{amsmath}
\usepackage{booktabs}
\usepackage{graphicx}
\usepackage{longtable}
\usepackage{listings}
\usepackage{xcolor}
\usepackage{hyperref}

\hypersetup{
  colorlinks=true,
  linkcolor=blue,
  urlcolor=blue,
  citecolor=blue
}
\lstset{
  basicstyle=\ttfamily\small,
  breaklines=true,
  columns=fullflexible,
  keepspaces=true,
  showstringspaces=false
}
\title{Windows Computer Literacy\\Setup and Everyday Reference}
\author{Hugh Rodgers}
\begin{document}
\maketitle

\LARGE{IF YOU FIND ANY COMMANDS OR INFORMATION ARE INCORRECT PLEASE LET ME KNOW, 123340696@umail.ucc.ie, I DONT USE WINDOWS SO I CANT VERIFY EVERYTHING}
\tableofcontents

\section{Install VS Code or VSCodium}
  Choose one editor. 
  The VS Code \href{https://code.visualstudio.com/docs/setup/windows}{User Setup} or VSCodium \href{https://vscodium.com/#install}{official installation options}; 
  Or use WinGet to install either editor.
  Winget install vscode --source winget.
  Winget install vscodium --source winget.

\section{Install the shared tools}
  \begin{itemize}

    \item Install Python following the instructions at \href{https://www.python.org/downloads/windows/}{official Windows downloads} 
    or with PowerShell using "winget install Python.Python.3.14".

    \item Install Git following the instructions at \href{https://git-scm.com/download/win}{the official Git downloads} 
    or with PowerShell using "winget install --id Git.Git -e --source winget".

    \item I have found many people have trouble with LaTeX in vscode on windows, i would still recomend this as it keeps all your things in one place.
    you can try \href{https://miktex.org/howto/install-miktex}{MiKTeX} or \href{https://tug.org/texlive/}{TeX Live}. 
    If none of this works and you dont care to install a local TeX distribution, use Overleaf is a browser alternative, this will cause problems down the line when you have longer document and it begins to ask for money.
    Other dedicated programs are available as well but i have no experience with them.
  \end{itemize}

  After installation you should verify that the tools are correctly installed by checking their versions.
  Close and reopen PowerShell after installation and type the following commands to check the versions of the installed tools.

  \begin{lstlisting}
    py --version
    git --version
    Get-Location
  \end{lstlisting}
  If \texttt{py} is unavailable but \texttt{python --version} identifies the installed interpreter, use \texttt{python} in the environment commands below.
  You may have to write python3 instead of python depending on the version
\newpage


\section{PowerShell, projects, and paths}
This is for PowerShell (Windows Terminal), not Command Prompt,
Git Bash, or WSL, for these instructions. A terminal runs a shell to navigate
folders and start programs. Relative paths begin from the current working
directory. In File Explorer, enable filename extensions so a source file does
not accidentally become \texttt{report.tex.txt}. Keep projects in a writable
user folder.

\begin{lstlisting}
pendulum-lab/
  README.md
  requirements.txt
  data/
    raw_measurements.csv
  analysis/
    analysis.py
    analysis.ipynb
  figures/
  report/
    report.tex
\end{lstlisting}

\begin{longtable}{p{0.50\linewidth}p{0.64\linewidth}}
\toprule Task & PowerShell command\\ \midrule
Current directory & \texttt{Get-Location} (alias: \texttt{pwd})\\
List files, including hidden & \texttt{Get-ChildItem -Force}\\
Move to the different directory & \texttt{Set-Location "DirectoryFilePath"}\\
Enter a path with spaces & \texttt{Set-Location "My Project"}\\
Copy a file & \texttt{Copy-Item source.csv copy.csv}\\
Move a file & \texttt{Move-Item copy.csv data.csv}\\
Find CSV files & \texttt{Get-ChildItem -Filter *.csv -Recurse}\\
Search file text & \texttt{Select-String -Path notes.txt -Pattern time}\\
Open the current folder & \texttt{Start-Process .}\\
Open the folder in the editor & \texttt{"code ."} or \texttt{"codium ."}\\
\bottomrule
\end{longtable}
Inspect paths before copying or moving files: an existing destination may be
overwritten. Use File Explorer's Recycle Bin when removing files. If the
editor command is unavailable, use File > Open Folder.
For more on PowerShell's file and directory commands, see Microsoft's
\href{https://learn.microsoft.com/en-us/powershell/}{working
with files and folders guide}.

\section{Python and project environments}
From the project directory, create a virtual environment and install packages using its interpreter:
\begin{lstlisting}
py -m venv .venv
.\.venv\Scripts\python.exe -m pip install numpy pandas matplotlib etc
.\.venv\Scripts\python.exe -m pip check
.\.venv\Scripts\python.exe -c "import sys; print(sys.executable)"
\end{lstlisting}
The printed path should include \texttt{.venv}. Optional PowerShell
activation is:
\begin{lstlisting}
.\.venv\Scripts\Activate.ps1
\end{lstlisting}
After activation, use \texttt{python -m pip ...}; run \texttt{deactivate} to
leave the environment. If activation is blocked, call the environment's
Python executable directly rather than weakening the execution policy.
A virtual environment keeps project packages separate; it is not a security
sandbox, and it should not be committed or copied to another computer. See the
\href{https://docs.python.org/3/library/venv.html}{Python virtual-environment
documentation}.

\section{Git and LaTeX}
Git works locally; GitHub hosts online copies and collaboration. Saving a file
is not the same as committing it, and a commit is not the same as pushing it.
For the shared commands and safe workflow, see the
\href{../Git/Git.pdf}{Git and GitHub guide}.

LaTeX source (\texttt{.tex}) is compiled into a PDF. Use Overleaf in a browser
or install MiKTeX/TeX Live for local builds. From the repository root, build
this handout with:
\begin{lstlisting}
latexmk -cd -pdf -interaction=nonstopmode -halt-on-error Windows/Windows.TEX
\end{lstlisting}
The shared \href{../LaTeX/LaTeX.pdf}{LaTeX guide} covers report structure,
equations, figures, and references; see also the
\href{https://www.overleaf.com/learn/latex/Learn_LaTeX_in_30_minutes}{Overleaf
LaTeX introduction} and
\href{https://www.latex-project.org/help/documentation/}{LaTeX Project
documentation}.

\section{Shortcuts and PowerShell profile}
\begin{itemize}
\item \texttt{Ctrl+C} interrupts a command; \texttt{Tab} completes paths;
\texttt{Up} recalls command history.
\item Windows Terminal commonly uses \texttt{Ctrl+Shift+C/V} for copy/paste.
\item \texttt{Win+E} opens File Explorer; \texttt{Win+Left/Right} tiles a
window to the side.
\end{itemize}

PowerShell reads a profile script when it starts. Display its location with
\verb|$PROFILE|. Create the parent folder and file only if needed, then open
the profile in Notepad:
\begin{lstlisting}
New-Item -ItemType Directory -Path (Split-Path -Parent $PROFILE) -Force | Out-Null
if (!(Test-Path -Path $PROFILE)) {
  New-Item -ItemType File -Path $PROFILE | Out-Null
}
notepad $PROFILE
\end{lstlisting}
Add personal shortcuts, save, and load them into the current session with
\begin{lstlisting}
function ll { Get-ChildItem -Force }
function croot { Set-Location $HOME }
. $PROFILE
\end{lstlisting}
Back up and review your profile before editing it. Bash aliases from
\texttt{.bashrc} do not automatically work in PowerShell. See Microsoft's
\href{https://learn.microsoft.com/en-us/powershell/module/microsoft.powershell.core/about/about_profiles}{PowerShell
profile documentation}.


\noindent For Python, fitting, and data-analysis exercises, use the
\href{../Example_Code/Exersize.ipynb}{shared exercise notebook}.
\end{document}
